\begin{abstract}
Robot learning is splitting into two bets: policies that bake competence into \emph{frozen
weights}---vision-language-action (VLA) models---and agents that write and refine their own
\emph{executable skills} as code. This survey organises the field around that axis,
\textbf{weights or skills}, and makes its central contribution a deep-dive that arranges
code-as-policy methods by their \emph{degree of self-improvement}: from zero-shot program
synthesis, through closed-loop self-repair and persistent skill memory, to the near-empty
cell where execution feedback, skill memory, and evolutionary search combine into one
open-ended loop. We anchor that frontier on ASPIRE and its real-world sibling ENPIRE, and we
map the complementary ``skills'' pole---from unsupervised reinforcement-learning skill
discovery to large-language-model skill libraries---showing that skills come in an RL flavour
and a code flavour, of which only the code flavour self-improves without gradients. Finally,
we connect the taxonomy to the emerging \emph{skill economy}: commercial robot-skill
marketplaces now distribute one-tap skills across robots, but ship only static playback, and
we lay out the open problems---adaptation, cross-embodiment portability, provenance, safety
verification, composition, and standardisation---that separate a store of animations from a
store of capabilities. The survey covers roughly seventy-five systems across six technique
families and provides a canonical taxonomy figure for each.
\end{abstract}
